\documentclass[11pt]{article}
\usepackage{amsmath,amsthm,hyperref}
\newtheorem{theorem}{Theorem}
\title{How many availability checks does a guaranteed rescue need?}
\author{}
\date{26 November 2026}
\begin{document}
\maketitle

A rescue controller can issue one action. Each action uses one of the offered sites $Q$, such as a warehouse berth or a sender's remaining transmission allowance. An old operation can hold a site while its final duty is still unresolved. The source contract gives a nonempty set of possible worlds, each with a set $W$ of sites the rescue can in fact use, all relative to its fixed demand. Availability does not change between the first inspection and commitment. This promise must come from the source or an authorized hold; a prediction of a rebind does not supply it. A selective inspection truthfully says whether one named site is usable, causes no change in reservations or availability, and can be followed immediately by commitment. There is a deadline expressed by a budget of such inspections, not by an assumed timing benchmark. No informative observation arrives otherwise. The controller must guarantee service in every admitted world, not maximize average rescue rate.

There are two admission laws. A verified-use rule requires a positive inspection of the chosen site. An inference-use rule may instead choose an uninspected site when all worlds consistent with the inspection transcript make it safe. This difference is about whether a proof from the source promise is acceptable evidence, not whether the site is physically available.

\begin{theorem}
Let $k(Q)$ be the smallest size of a subset of $Q$ meeting every possible $W$, or infinity if none exists. If $k(Q)$ is finite, the exact worst-case number of inspections is $k(Q)$ under verified use and $k(Q)-1$ under inference use. With nonnegative deterministic costs $c_x$ and a worst-case budget, replace these bounds by the minimum over meeting sets $K\subseteq Q$ of respectively
\[
  \sum_{x\in K}c_x
  \qquad\text{and}\qquad
  \sum_{x\in K}c_x-\max_{x\in K}c_x.
\]
The canonical summary for all future offers and budgets is the antichain of inclusion-minimal possible $W$. In the worst case, at least $\binom{n}{\lfloor n/2\rfloor}$ bits are needed in a fixed-width exact summary for $n\ge2$ named sites.
\end{theorem}

\begin{proof}
Before the first positive reply the controller sees only a sequence of negatives; along this branch adaptive selection cannot create a second path except by obtaining a positive, which already gives a safe action. Once there has been a positive it cannot improve the worst-case inspection cost to keep searching. Remove repeated queries. Truncate the all-negative list when the next query is guaranteed positive, or, for inference use, when the unqueried final choice is guaranteed safe. Every world must contain a positive site in this list of queried candidates plus the possible final inferred site, so the list meets every world. Conversely, inspect a meeting set in any order and commit on the first positive. If every earlier result is negative and inference is allowed, the last site is necessarily usable and need not be inspected. For the cost claim use an inclusion-minimal meeting set: for every member there is a world in which that member alone from the set is positive. Hence the final position in the list really determines the worst-case sum. With inference put a largest-cost item last. In the lower-bound direction, a minimum-cost algorithm can likewise stop the negative branch as soon as no compatible world would return negative at its next inspection; its paid list is itself a meeting set, except for at most the terminal inferred choice. Thus no impossible branch is charged. Randomization cannot improve a sure result with a hard budget.

A set meets every $W$ iff it meets every minimal $W$. Conversely the offers admitting no service even with unrestricted inspection are precisely the subsets of complements of some $W$. Their maximal elements give the complements of the minimal sets. Thus two different minimal antichains are separated by an offer alone. All nonempty families chosen from the middle level of the Boolean lattice are such antichains, giving $2^{\binom{n}{\lfloor n/2\rfloor}}-1$ distinct classes and the stated bit bound. The consumer drops the occurrence law inside each non-minimal world, the identities of old attempts that caused occupancy, and anything about changing sites before commitment. Those facts can matter to a different consumer.
\end{proof}

\paragraph{Bounded false negatives.}
Suppose a reported positive remains conclusive but a free-site inspection can return negative up to $b$ times across the whole attempt. The controller cannot distinguish a false negative from a real negative. Repeat inspections are allowed. The same argument gives an exact bounded-cover problem. Let $m_x$ be the number of inspections of site $x$ on the no-positive branch. For verified use the requirement is $\sum_{x\in W\cap Q}m_x\ge b+1$ for every world; for inference use with terminal choice $y\in Q$ it is required only in worlds with $y\notin W$. Minimize $\sum_x c_xm_x$, also choosing $y$ in the inference case. A false-negative budget can be spent across sites, so repeating only a narrow set or probing more sites are competing options.

Indeed, an all-negative record can occur in exactly those worlds with at most $b$ hits among the scheduled inspections. If a positive arrives the controller can immediately commit. If none arrives, the verified policy must have ruled that branch out and the inference policy must have ruled out every remaining world where $y$ is not usable. This proves both directions, including arbitrary adaptive policies by following their single no-positive branch. The minimal antichain still suffices because every count is nonnegative and the inference condition is only imposed on sets missing $y$. The case $b=0$ recovers the original cost formula. The replay branches directly on physically supported positives and on negatives that consume a possible false-negative credit. False positives, recovery work induced by the read, and a failure quota that refreshes with each supplier are different models.

\paragraph{Same separate histories, different cost.}
For $n\ge3$, one possible law has every singleton world with probability $1/[n(n-1)]$ and the all-free world with probability $(n-2)/(n-1)$. Another has every set with exactly one unavailable site with probability $1/n$. Individual sites have the same free probability $(n-1)/n$ under both; if the whole relevant availability interval is constant, each individual time profile has the same law too. Neither world family admits a blind choice, and both support guaranteed rescue. But the first requires $n-1$ inspections under inference use and the second requires one. The joint support, not per-site logging accuracy, determines the budget. At zero extra inspection budget both controllers refuse rather than guessing. Checking whether $k(Q)$ exceeds a budget is exactly the ordinary hitting-set problem, not a new complexity claim. The replay compares the rule to an independent adaptive decision tree, including offers that cannot work.

\paragraph{Relation to checked sources.}
In the pinned two-site RWARE replay an admitted old forward move reaches and holds its staged site; in the SubT replay the old staged send consumes that sender's rolling allowance and makes native progress while final delivery remains owed. The new replay reads safe rescue actions from those two existing source adapters under held scopes, not from a declared answer. Their verified two-site cases project to the same singleton availability sets and admit the same exact probe bound. The prior rebind contrast is a negative control, not evidence that a scope remains stable. These fixture results do not give an actual inspection latency, a stock read API, physical berth clearance outside the model, or an independent source selection. The normalized larger antichains and equal-marginal laws are abstract workloads, not scalability results for either native adapter.

The strongest matched comparator is a direct contingent policy over exactly the same worlds, observations and permissions. It agrees. The smaller reusable antichain and the distinction between a proof-only final site and a required fresh positive are the returned results. They can inform a receiver interface but do not establish an operational gain over the matched comparator or discharge the integration test.
\end{document}
